\chapter{Arquitectura del sistema}

\section{Visión general}

La fase 1 se implementa como el paquete Python \texttt{decisionlab}. Su objetivo
es transformar problemas en lenguaje natural en modelos de decisión ejecutables.
La arquitectura separa tres capas:

\begin{itemize}
\item Capa de agentes: investigación, formalización, razonamiento y construcción.
\item Capa de artefactos: informes, formulaciones, especificaciones, código,
tests, trazas y estado.
\item Capa de memoria: grafo, recuperación densa/dispersa y ciclo de vida.
\end{itemize}

\begin{figure}[H]
\centering
\begin{verbatim}
Usuario / CLI / Web
        |
        v
   +---------+
   | Router  |
   +----+----+
        |
  +-----+------+------+
  |            |      |
  v            v      v
Agentes   Artefactos Memoria
R/F/R/B   MinIO+PG   Neo4j
          codigo     Qdrant
                     PG lifecycle
\end{verbatim}
\caption{Arquitectura lógica de la fase 1.}
\end{figure}

\section{Router como máquina de estados}

El Router es el coordinador principal. No delega el flujo global en un modelo de
lenguaje, sino que mantiene una máquina de estados explícita. Esto permite
reanudar ejecuciones, registrar trazas, pedir revisión humana y decidir cuándo
escribir en memoria.

\begin{figure}[H]
\centering
\begin{verbatim}
CLASSIFY_UMBRELLA
  -> RESEARCH -> REVIEW_RESEARCH -> MEMORY_RESEARCH
  -> FORMALIZE -> REVIEW_FORMALIZE -> MEMORY_FORMALIZE
  -> GET_ENV_SPEC
  -> REASON -> REVIEW_REASON -> MEMORY_REASON
  -> BUILD -> REVIEW_BUILD -> MEMORY_BUILD
  -> DONE
\end{verbatim}
\caption{Máquina de estados simplificada del pipeline.}
\end{figure}

Las etapas \texttt{MEMORY\_*} se ejecutan solo si la infraestructura de memoria
está disponible. Además, se colocan después de revisión humana. Esta decisión
evita que borradores rechazados o especificaciones inválidas contaminen el grafo
de conocimiento.

\section{Separación de responsabilidades}

Cada agente tiene una responsabilidad estrecha. Esta separación reduce el tamaño
cognitivo de cada prompt y facilita repetir una etapa sin repetir todo el
pipeline.

\begin{table}[H]
\begin{center}
\begin{tabular}{|l|p{9cm}|} \hline
\textbf{Componente} & \textbf{Responsabilidad} \\ \hline
Router & Coordina estados, revisión, persistencia, memoria y reejecuciones. \\ \hline
Researcher & Identifica paradigmas científicos y produce informes. \\ \hline
Formalizer & Convierte informes en formulaciones matemáticas comparables. \\ \hline
Reasoner & Adapta formulaciones al entorno y genera especificaciones JSON. \\ \hline
Builder & Escribe modelos Python y pruebas ejecutables. \\ \hline
MemoryAgent & Extrae hechos, entidades y relaciones de artefactos aceptados. \\ \hline
\end{tabular}
\caption{Responsabilidades principales.}
\end{center}
\end{table}

\section{Artefactos}

La ejecución produce artefactos persistentes. MinIO almacena los bytes y
PostgreSQL almacena metadatos. Esta separación evita guardar informes largos o
código en tablas relacionales, pero mantiene la capacidad de listar, auditar y
registrar resultados.

\begin{verbatim}
research/{run_id}/
  report.md
  deep/{paradigm}.md
  formulations/{paradigm}.md
  env_spec.json
  pipeline_state.json
  trace.jsonl

models/{run_id}/
  reasoner/{paradigm}/{formulation}.json
  builder/{paradigm}/{formulation}_model.py
  builder/{paradigm}/test_{formulation}.py
  builder/{paradigm}/{formulation}_validation.json
\end{verbatim}

\section{Infraestructura compartida}

La fase 1 no crea clientes globales de base de datos, almacenamiento o memoria.
Recibe un contenedor de servicios inicializado por el punto de entrada:

\begin{lstlisting}[language=Python]
Services(
    db=DatabaseService,
    storage=StorageService,
    kg=KnowledgeGraph | None,
    vectors=VectorStore | None,
    embeddings=EmbeddingService | None,
)
\end{lstlisting}

PostgreSQL y MinIO son obligatorios para una ejecución completa. Neo4j, Qdrant,
embeddings y reranking son opcionales: si faltan, el pipeline sigue funcionando
en modo degradado, aunque pierde recuperación y memoria.

\section{Contrato con la fase 2}

El producto final no depende de clases internas de la fase 2. Los modelos
generados implementan un contrato por duck typing:

\begin{lstlisting}[language=Python]
def decide(self, perception: dict) -> Action
def update(self, action, reward, new_perception) -> None
def get_state(self) -> dict
\end{lstlisting}

Esto permite cargar modelos dinámicamente en el simulador. La regla principal es
que \texttt{decide} sea de solo lectura y que toda mutación del aprendizaje se
realice en \texttt{update}.
